\section{System requirements and problem analysis}

\subsection{Problem context}
Ecopark là khu đô thị xanh có nhiều điểm tham quan, hồ nước và công viên. Nhu cầu
di chuyển ngắn trong nội khu phù hợp với xe đạp, nhưng nếu quản lý bằng giấy tờ
hoặc bảng tính rời rạc thì nhân sự vận hành khó biết chính xác xe nào đang ở bãi
nào, xe nào đang được thuê, xe nào hỏng, khách nào đủ điều kiện thuê và doanh thu
theo từng kỳ là bao nhiêu.

Bài toán đặt ra là xây dựng một hệ thống thông tin hỗ trợ thuê-trả xe đạp end-to-end
cho hai nhóm người dùng chính: khách thuê xe và nhân sự vận hành. Hệ thống cần
đáp ứng cả trải nghiệm khách hàng, kiểm soát nghiệp vụ tại bãi và báo cáo quản lý
tập trung.

\subsection{Problems to solve}
\begin{itemize}
  \item \textbf{Dữ liệu xe phân tán:} mỗi điểm tập kết có nhiều xe, xe có mã duy
  nhất và trạng thái thay đổi liên tục. Nếu cập nhật thủ công, người quản lý khó
  tìm nhanh vị trí/trạng thái thật của từng xe.
  \item \textbf{Luồng thuê cần kiểm soát:} khách phải có tài khoản, email, số điện
  thoại, địa chỉ và giấy tờ định danh; khi nhận xe cần đặt cọc 200k và để lại
  CCCD hoặc giấy tờ tùy thân.
  \item \textbf{Trả xe linh hoạt:} khách có thể trả xe ở bất kỳ bãi Ecopark nào,
  nên hệ thống phải cập nhật lại vị trí xe, trạng thái xe và vé thuê ngay tại
  thời điểm trả.
  \item \textbf{Tính phí có nhiều nhánh:} giá thuê phụ thuộc gói giờ, phụ thu trễ
  tính theo block 30 phút và cư dân Ecopark hợp lệ được giảm 40\% phí thuê cơ bản.
  \item \textbf{Quản trị cần báo cáo:} Admin/Operator cần thống kê theo ngày,
  tuần, tháng, bãi xe và xe để theo dõi vận hành thay vì chỉ xem từng lượt thuê.
\end{itemize}

\subsection{Functional requirements}
\begin{longtable}{P{0.13\textwidth}P{0.32\textwidth}P{0.45\textwidth}}
\toprule
\textbf{Mã} & \textbf{Requirement} & \textbf{Mô tả nghiệm thu} \\
\midrule
FR-01 & Account and identity management & Khách tạo tài khoản lần đầu với email, số điện thoại, địa chỉ, CCCD/CMND và mật khẩu; hệ thống kiểm tra định dạng, dữ liệu trùng, độ mạnh mật khẩu, email verification demo-local và giấy tờ bị khóa. \\
FR-02 & Resident discount & Khách cư dân khai báo địa chỉ/thẻ cư dân; hệ thống lưu trạng thái đủ điều kiện giảm 40\%. \\
FR-03 & Station search & Khách tìm bãi gần mình theo GPS demo, vị trí đã lưu hoặc vị trí nhập tay, chọn phạm vi và xem số xe rảnh. \\
FR-04 & Bike selection and request & Khách lọc loại xe, chọn xe đang rảnh và gửi yêu cầu thuê với thời lượng 1/2/3 giờ; hệ thống chỉ tạo request khi tọa độ khách nằm trong bán kính bãi nhận. \\
FR-05 & Handover and deposit & Staff/Admin kiểm tra request theo scope bãi, đối chiếu giấy tờ, xác nhận cọc 200k, đổi xe trước handover khi cần, kiểm tra lại geofence khách ở bãi nhận, giữ giấy tờ và chuyển xe sang đang thuê. \\
FR-06 & Return and ticket & Customer xác nhận bãi trả bằng GPS trong bán kính bãi; Staff nhận xe trong bãi được phân công, Admin/Operator xử lý toàn hệ thống, hỗ trợ trả khác bãi, kiểm tra tình trạng xe, tính phí, phụ thu, giảm giá và xuất vé. \\
FR-07 & Fleet management & Admin/Operator quản lý bãi, xe, trạng thái xe, vị trí xe và lý do đổi trạng thái. \\
FR-08 & Reports and audit dashboard & Admin/Operator xem metric, biểu đồ doanh thu/trạng thái đội xe/lượt thuê-phụ thu, lọc theo bãi/xe, xuất CSV và xem audit log thay đổi trạng thái xe. \\
FR-09 & Concurrent demo sessions & Hệ thống cho nhiều browser/context đăng nhập độc lập để demo khách, admin và tab GPS cùng lúc. \\
\bottomrule
\end{longtable}

\subsection{Non-functional requirements}
\begin{longtable}{P{0.20\textwidth}P{0.70\textwidth}}
\toprule
\textbf{Nhóm yêu cầu} & \textbf{Nội dung} \\
\midrule
Usability & Giao diện phải đọc được trên desktop/mobile, tránh tràn chữ, giữ flow thuê xe liền mạch và phân biệt rõ khách hàng với nhân sự vận hành. \\
Reliability & Trạng thái xe không được bị double-booking; xe đã giữ chỗ hoặc đang thuê không được xuất hiện như xe có thể thuê. \\
Maintainability & Code tách rõ front-end, API, schema, pricing và test; tài liệu PDF/README phản ánh đúng hành vi triển khai. \\
Performance & Chạy mượt trong môi trường demo local và khi phục vụ qua GCP VM/container; animation ưu tiên transform/opacity và có hỗ trợ \texttt{prefers-reduced-motion}. \\
Security baseline & Có role guard, staff scope theo bãi, session cookie HttpOnly, PBKDF2 salted password hash, token verify/reset lưu dạng hash và kiểm tra input server-side; secret production được định hướng đặt trong GCP Secret Manager. \\
Portability & Có thể chạy trên một máy bằng Node.js 24 và SQLite để demo, đồng thời deploy lên GCP Compute Engine với persistent disk cho file SQLite. \\
\bottomrule
\end{longtable}

\subsection{Submission checklist mapping}
\begin{longtable}{P{0.40\textwidth}P{0.50\textwidth}}
\toprule
\textbf{Mục yêu cầu trong đề} & \textbf{Cách report này đáp ứng} \\
\midrule
System requirements and problem analysis & Section 1 phân tích bài toán, requirement chức năng/phi chức năng và phạm vi. \\
Use-case analysis & Section 2 mô tả actor, quy tắc nghiệp vụ, UC001--UC005 và sơ đồ use case. \\
User stories & Section 3 chuyển requirement thành story theo vai trò Customer, Staff và Admin. \\
System architecture & Section 4 mô tả runtime architecture, luồng dữ liệu và đánh giá back-end. \\
Component architecture and design & Section 5 mô tả component front-end/API/domain/database và lifecycle thuê xe. \\
Database design & Section 6 trình bày relational schema, quan hệ bảng và ràng buộc nghiệp vụ. \\
User interface design & Section 7 trình bày phong cách UI, màn hình chính và screenshot. \\
Component implementation & Section 8 ánh xạ component/code module với use case. \\
API design and implementation & Section 9 liệt kê endpoint, dữ liệu vào/ra và quy tắc xử lý. \\
Testing and evaluation & Section 10 tổng hợp test backend, smoke UI/UC, kết quả và giới hạn. \\
\bottomrule
\end{longtable}
